% Условие убывания: направление p_k годится, если оно образует тупой угол
% с градиентом, то есть лежит в полуплоскости grad^T f . p < 0.
\begin{figure}[htbp]
\centering
\begin{tikzpicture}[>={Stealth[length=2.2mm]}, scale=1.0]
  % линии уровня
  \foreach \r in {0.55, 1.0, 1.45, 1.9} {
    \draw[blue!35!black, rotate=22] (0,0) ellipse ({\r*1.6} and {\r*0.95});
  }

  \coordinate (XK) at (-1.55,-0.95);
  \fill (XK) circle (2.0pt);
  \node[font=\footnotesize, anchor=north east] at (XK) {$x_k$};

  % полуплоскость направлений убывания
  \begin{scope}
    \clip (XK) circle (2.15);
    \fill[teal!12, rotate around={-32:(XK)}]
      ($(XK)+(0,-2.4)$) rectangle ($(XK)+(2.6,2.4)$);
  \end{scope}
  % касательная к линии уровня (граница полуплоскости)
  \draw[dashed, gray, rotate around={-32:(XK)}]
    ($(XK)+(0,-2.0)$) -- ($(XK)+(0,2.0)$);

  % градиент и антиградиент
  \draw[->, red!70!black, very thick] (XK) -- ++(-1.05,-0.66);
  \node[font=\footnotesize, red!70!black, anchor=north east]
    at ($(XK)+(-1.0,-0.62)$) {$\nabla f(x_k)$};
  \draw[->, teal!65!black, very thick] (XK) -- ++(1.05,0.66);
  \node[font=\footnotesize, teal!65!black, anchor=south]
    at ($(XK)+(1.0,0.68)$) {$-\nabla f(x_k)$};

  % произвольное допустимое направление
  \draw[->, orange!85!black, very thick] (XK) -- ++(1.5,-0.25);
  \node[font=\footnotesize, orange!80!black, anchor=west]
    at ($(XK)+(1.55,-0.25)$) {$p_k$};

  \node[font=\scriptsize, teal!50!black, anchor=west] at (0.45,-1.95)
    {$\nabla\T f(x_k)\, p_k < 0$};
\end{tikzpicture}
\caption{Направление $p_k$ годится как направление спуска ровно тогда, когда
оно лежит в открытой полуплоскости $\nabla\T f(x_k)\, p_k < 0$ (закрашена).
Антиградиент — «самое глубокое» направление в ней, его и берут в градиентных
методах первого порядка.}
\label{fig:descent-cone}
\end{figure}
